% CVPR 2026-based Yuvion technical report template.
% Official CVPR author kit: https://github.com/cvpr-org/author-kit

\documentclass[10pt,twocolumn,letterpaper]{article}

% Yuvion public technical-report mode.
\usepackage[pagenumbers]{cvpr}

\usepackage[utf8]{inputenc}
\usepackage[T1]{fontenc}
\usepackage{microtype}
\usepackage{graphicx}
\usepackage{booktabs}
\usepackage{amsmath}
\usepackage{amssymb}
\usepackage{multirow}
\usepackage{algorithm}
\usepackage{algorithmic}
\usepackage{fancyhdr}
\usepackage{xcolor}
\usepackage[pagebackref,breaklinks,colorlinks,allcolors=cvprblue]{hyperref}

% Yuvion technical-report header layered on top of the CVPR layout.
\newlength{\yuvionheadextra}
\setlength{\yuvionheadextra}{24pt}
\addtolength{\yuvionheadextra}{-\headheight}
\addtolength{\topmargin}{-\yuvionheadextra}
\setlength{\headheight}{24pt}
% Move the branded header into the existing top margin without moving the body.
\addtolength{\headsep}{4pt}
\addtolength{\topmargin}{-4pt}
\fancypagestyle{yuvionheader}{%
  \fancyhf{}%
\fancyhead[L]{\raisebox{0pt}{\includegraphics[height=17pt]{alibaba-security-en.png}}}%
  \fancyhead[C]{}%
\fancyhead[R]{\raisebox{0pt}{\includegraphics[height=17pt]{yuvion_en_good.png}}}%
  \fancyfoot[C]{\thepage}%
  \renewcommand{\headrulewidth}{0.8pt}%
  \renewcommand{\footrulewidth}{0pt}%
}

\def\paperID{}
\def\confName{CVPR}
\def\confYear{2026}

\title{Yuvion Team Paper Template and Author Guidelines}

\author{First Author\\
Yuvion Team, Alibaba Group\\
{\tt\small first.author@example.com}
\and
Second Author\\
Yuvion Team, Alibaba Group\\
{\tt\small second.author@example.com}
}

\begin{document}
\maketitle
\pagestyle{yuvionheader}
\thispagestyle{yuvionheader}

\begin{abstract}
The ABSTRACT is to be in fully justified italicized text, at the top of the left-hand column, below the author and affiliation information.
Use the word ``Abstract'' as the title, in 12-point Times, boldface type, centered relative to the column, initially capitalized.
The abstract is to be in 10-point, single-spaced type.
Leave two blank lines after the Abstract, then begin the main text.
Look at previous \confName abstracts to get a feel for style and length.
\end{abstract}

\section{Using This Template}
\label{sec:intro}

This CVPR-based template is maintained for papers authored by the Yuvion Team at Alibaba Group. It preserves the standard CVPR two-column typography and common formatting conventions while adding the Yuvion-branded page header.

\subsection{Affiliation}

Use \textbf{Yuvion Team, Alibaba Group} as the affiliation for team authors. Replace the placeholder author names and email addresses in the main \texttt{.tex} file before releasing the paper.

\subsection{Language}

All manuscripts should be written in English.

\subsection{Mathematics}

Please, number all of your sections and displayed equations as in these examples:
\begin{equation}
  E = m\cdot c^2
  \label{eq:important}
\end{equation}
and
\begin{equation}
  v = a\cdot t.
  \label{eq:also-important}
\end{equation}
It is important for the reader to be able to refer to any particular equation.
Just because you did not refer to it in the text does not mean that some future reader might not need to refer to it.
It is cumbersome to have to use circumlocutions like ``the equation second from the top of page 3 column 1''.
All authors will benefit from reading Mermin's description of how to write mathematics:
\url{http://www.pamitc.org/documents/mermin.pdf}.

\subsection{Miscellaneous}

\noindent
Compare the following:\\
\begin{tabular}{ll}
 \verb'$conf_a$' &  $conf_a$ \\
 \verb'$\mathit{conf}_a$' & $\mathit{conf}_a$
\end{tabular}\\
See The \TeX book, p165.

The space after \eg, meaning ``for example'', should not be a sentence-ending space.
So \eg is correct, {\em e.g.} is not.
The provided \verb'\eg' macro takes care of this.

When citing a multi-author paper, you may save space by using ``et alia'', shortened to ``\etal'' (not ``{\em et.\ al.}'' as ``{\em et}'' is a complete word).
If you use the \verb'\etal' macro provided, then you need not worry about double periods when used at the end of a sentence as in Alpher \etal.
However, use it only when there are three or more authors.
Thus, the following is correct:
   ``Frobnication has been trendy lately.
   It was introduced by Alpher~\cite{Alpher02}, and subsequently developed by
   Alpher and Fotheringham-Smythe~\cite{Alpher03}, and Alpher \etal~\cite{Alpher04}.''

This is incorrect: ``... subsequently developed by Alpher \etal~\cite{Alpher03} ...'' because reference~\cite{Alpher03} has only two authors.

\section{Paper Publication and Open-Source Workflow}

\subsection{Paper publication approval}

Before submitting or publicly releasing a paper, complete the internal publication approval process through the \href{https://yida-group.alibaba-inc.com/APP_ZHVQTG6LH0I0BOSUJDQ8/workbench/FORM-82EC232FEEE9411BB171ECBA40FCD614T9QF?spm=ata.25287382.0.0.27147536laDwbn}{Yida paper publication approval form}. Complete the required review and approval steps before any external publication.

\subsection{Code open-source approval}

If the project will release source code, models, datasets, or related technical assets, initiate and complete the compliance review through the \href{https://lc.alibaba-inc.com/\#/lc/home?survey=1}{Alibaba open-source approval platform} before public release.

\section{Formatting your paper}
\label{sec:formatting}

All text must be in two-column format.
The total allowable size of the text area is $6\frac78$ inches (17.46 cm) wide by $8\frac78$ inches (22.54 cm) high.
The columns should be $3\frac14$ inches (8.25 cm) wide, with a $\frac{5}{16}$ inch (0.8 cm) space between them.
The main title (on the first page) should begin 1 inch (2.54 cm) from the top edge of the page.
The second and following pages should begin 1 inch (2.54 cm) from the top edge.
On all pages, the bottom margin should be $1\frac{1}{8}$ inches (2.86 cm) from the bottom edge of the page for $8.5 \times 11$-inch paper;
for A4 paper, approximately $1\frac{5}{8}$ inches (4.13 cm) from the bottom edge of the
page.

%-------------------------------------------------------------------------
\subsection{Margins and page numbering}

All printed material, including text, illustrations, and charts, must be kept
within a print area $6\frac{7}{8}$ inches (17.46 cm) wide by $8\frac{7}{8}$ inches (22.54 cm)
high.
%
Page numbers should be in the footer, centered, and $\frac{3}{4}$ inches from the bottom of the page.
This template places page numbers in the centered footer.



%-------------------------------------------------------------------------
\subsection{Type style and fonts}

Wherever Times is specified, Times Roman may also be used.
If neither is available on your word processor, please use the font closest in
appearance to Times to which you have access.

MAIN TITLE.
Center the title $1\frac{3}{8}$ inches (3.49 cm) from the top edge of the first page.
The title should be in Times 14-point, boldface type.
Capitalize the first letter of nouns, pronouns, verbs, adjectives, and adverbs;
do not capitalize articles, coordinate conjunctions, or prepositions (unless the title begins with such a word).
Leave two blank lines after the title.

AUTHOR NAME(s) and AFFILIATION(s) are to be centered beneath the title
and printed in Times 12-point, non-boldface type.
This information is to be followed by two blank lines.

The ABSTRACT and MAIN TEXT are to be in a two-column format.

MAIN TEXT.
Type main text in 10-point Times, single-spaced.
Do NOT use double-spacing.
All paragraphs should be indented 1 pica (approx.~$\frac{1}{6}$ inch or 0.422 cm).
Make sure your text is fully justified---that is, flush left and flush right.
Please do not place any additional blank lines between paragraphs.

The captions of figures and tables should be in 9-point Roman type.
Short captions should be centered.  Table captions should be above tables, while figure captions should be below figures.

\noindent Callouts should be 9-point Helvetica, non-boldface type.
Initially capitalize only the first word of section titles and first-, second-, and third-order headings.

FIRST-ORDER HEADINGS.
(For example, {\large \bf 1. Introduction}) should be Times 12-point boldface, initially capitalized, flush left, with one blank line before and one blank line after.

SECOND-ORDER HEADINGS.
(For example, { \bf 1.1. Database elements}) should be Times 11-point boldface, initially capitalized, flush left, with one blank line before and one after.
If you require a third-order heading (we discourage it), use 10-point Times, boldface, initially capitalized, flush left, preceded by one blank line, followed by a period, and your text on the same line.

%-------------------------------------------------------------------------
\subsection{Footnotes}

Please use the footnotes\footnote{This is what a footnote looks like.
It often distracts the reader from the main flow of the argument.} sparingly.
Indeed, try to avoid footnotes altogether and include necessary peripheral observations in the text (within parentheses, if you prefer, as in this sentence).
If you wish to use a footnote, place it at the bottom of the column on the page on which it is referenced.
Use Times 8-point type, single-spaced.


%-------------------------------------------------------------------------
\subsection{Cross-references}

For the benefit of author(s) and readers, please use the
{\small\begin{verbatim}
  \cref{...}
\end{verbatim}}  command for cross-referencing to figures, tables, equations, or sections.
This will automatically insert the appropriate label alongside the cross-reference as in this example:
\begin{quotation}
  To see how our method outperforms previous work, see \cref{tab:example} and \cref{eq:important}.
  It is also possible to refer to multiple targets at once, \eg~to \cref{eq:important,eq:also-important}.
  You may also return to \cref{sec:formatting} or look at \cref{eq:also-important}.
\end{quotation}
If you do not wish to abbreviate the label, for example, at the beginning of the sentence, you can use
{\small\begin{verbatim}
  \Cref{...}
\end{verbatim}}
command. Here is an example:
\begin{quotation}
  \Cref{tab:example} is also quite important.
\end{quotation}

%-------------------------------------------------------------------------
\subsection{References}

List and number all bibliographical references in 9-point Times, single-spaced, at the end of your paper.
When referenced in the text, include the citation number in square brackets, for
example~\cite{Authors14}.
Where appropriate, include page numbers and the name(s) of editors of referenced books.
When citing multiple papers at once, make sure that you cite them in numerical order such as this \cite{Alpher02,Alpher03,Alpher05,Authors14b,Authors14}.
If you use the template as advised, this will be taken care of automatically.

\begin{table}
  \caption{Results.   Ours is better.}
  \label{tab:example}
  \centering
  \begin{tabular}{@{}lc@{}}
    \toprule
    Method & Frobnability \\
    \midrule
    Theirs & Frumpy \\
    Yours & Frobbly \\
    Ours & Makes one's heart Frob\\
    \bottomrule
  \end{tabular}
\end{table}

%-------------------------------------------------------------------------
\subsection{Illustrations, graphs, and photographs}

All graphics should be centered.
In \LaTeX, avoid using the \texttt{center} environment for this purpose, as this adds potentially unwanted whitespace.
Instead, use
{\small\begin{verbatim}
  \centering
\end{verbatim}}
at the beginning of your figure.
Please ensure that any point you wish to make is resolvable in a printed copy of the paper.
Resize fonts in figures to match the font in the body text, and choose line widths that render effectively in print.
Readers, even of an electronic copy, may choose to print the paper.
You cannot insist that they do otherwise and, therefore, must not assume that they can zoom in to see tiny details on a graphic.

When placing figures in \LaTeX, it is almost always best to use \verb+\includegraphics+, and to specify the figure width as a multiple of the line width as in the example below
{\small\begin{verbatim}
   \usepackage{graphicx} ...
   \includegraphics[width=0.8\linewidth]
                   {myfile.pdf}
\end{verbatim}
}


%-------------------------------------------------------------------------
\subsection{Color}

Please refer to the author guidelines on the \confName\ \confYear\ web page for a discussion of the use of color in your document.

If you use color in your plots, please keep in mind that a significant subset of readers may have a color vision deficiency; red-green blindness is the most frequent kind.
Hence, avoid relying only on color as the discriminative feature in plots (such as red \vs green lines), but add a second discriminative feature to ease disambiguation.

\section{Final Checklist}

Before external release, verify the author list, use \textbf{Yuvion Team, Alibaba Group} as the affiliation, complete the paper publication approval process, and complete the open-source review whenever code or related assets will be released.

{\small
\bibliographystyle{ieeenat_fullname}
\bibliography{yuvion_team_template}
}

\end{document}
